\documentclass[10pt,a4paper]{amsart}
\usepackage[margin=19mm]{geometry}
\usepackage{amsmath,amssymb,mathtools}
\usepackage[scr=rsfs,cal=euler]{mathalfa}
\usepackage{xfrac}
\usepackage[unicode,hidelinks,bookmarks=false]{hyperref}
\newcommand{\ie}{i.e.\ }
\newcommand{\cf}{cf.\ }
\newcommand{\resp}{respectively\ }
\newcommand{\Subsection}{\subsection}
\providecommand{\nobreakdash}{\nobreak\hbox{-}}
\setlength{\emergencystretch}{2em}
\multlinegap0pt
\usepackage{bm}
\usepackage{bbm}
\usepackage{dsfont}
\usepackage{xspace}
\usepackage{cancel}
\usepackage{mathtools}
\usepackage{amsthm}
\makeatletter
\@ifundefined{lemma}{\newtheorem{lemma}{Lemma}}{}
\@ifundefined{theorem}{\newtheorem{theorem}{Theorem}}{}
\makeatother
\title[Minimizing the Makespan Approximately on Two Identical Paral]{Minimizing the Makespan Approximately on Two Identical Parallel Machines with a Loading--Unloading Server}
\author{Keramat Hasani, Frank Werner}
\date{Source: arXiv:2608.18837v1 (CC BY 4.0)}
\begin{document}
\maketitle

\noindent\emph{Source and licence.} Minimizing the Makespan Approximately on Two Identical Parallel Machines with a Loading--Unloading Server. Version arXiv:2608.18837v1 (CC BY 4.0), distributed under the Creative Commons Attribution 4.0 International licence, as recorded in the arXiv licence metadata for this exact version and shown on the abstract page https://arxiv.org/abs/2608.18837v1. Full rights notice: licenses/arxiv-2608.18837-CC-BY-4.0.txt.\par\smallskip

\noindent\textit{Editorial scope.} Counted unit: the Theorem whose source label is \texttt{thm:lpt-unrestricted-upper} in the article (source0.tex lines 1265--1275, source label \texttt{thm:lpt-unrestricted-upper}), delivered together with the article's own complete original proof of it (source0.tex lines 1277--1433, one \texttt{proof} environment of 157 lines). No mathematics is written, repaired or re-derived: statement and proof are the article's own text line for line. Required context retained uncounted: eq:lpt-machine-lb (lines 1066--1069); eq:lpt-short-lb (lines 1071--1076); thm:lpt-long-jobs (lines 1090--1100); lem:lpt-prefix-tail (lines 1140--1168); eq:lpt-prefix-idle (lines 1159--1162); eq:lpt-prefix-tail-length (lines 1164--1167); eq:lpt-delete-short (lines 1205--1208). Every definition, display and label of those lines is retained unchanged. Omitted from this package and not redistributed: 1--1065, 1070--1070, 1077--1089, 1101--1139, 1163--1163, 1168--1204, 1209--1264, 1276--1276, 1434--2415. Nothing inside a retained excerpt is omitted or altered: every retained body is delivered exactly as the article's lines are written, so no omission receipt of any kind accompanies this package. Citations: the retained excerpts carry no citation command at all, so no bibliography entry of the article is transcribed and no reference is redistributed. Reference adaptation: none was needed and none was made; every reference inside the delivered unit resolves inside it, so no unresolved reference or citation is produced. Printed slips kept literal: no printed word, symbol, hypothesis or number of the counted statement or of its proof was altered. Layout and class: the article's own class is \texttt{article} with packages including geometry, amsmath, amssymb, amsthm, mathtools, natbib, hyperref; the wrapper is the lane's pinned \texttt{amsart} editorial wrapper at 10pt on A4, which restates the theorem-like environments the retained text needs and adds the article's own macro definitions that the retained text needs (none), with extra wrapper packages (bm, bbm, dsfont, xspace, cancel, mathtools). The delivered numbering is the unit's own. Assets: no retained span contains a figure environment, a graphics-inclusion command, a PDF-page inclusion, a table or a TikZ picture; no image file is copied into this package, the package declares no asset dependency and no image file is redistributed with it. Licence: version arXiv:2608.18837v1 of this article is distributed under the Creative Commons Attribution 4.0 International licence, as recorded in the arXiv licence metadata for this exact version and shown on the abstract page https://arxiv.org/abs/2608.18837v1. A full rights notice is delivered at licenses/arxiv-2608.18837-CC-BY-4.0.txt. Characters: every retained excerpt is pure ASCII, so the delivered engine, XeLaTeX with its default Latin Modern text font, reports no missing character.
\begin{equation}
C^*\ge \frac{E+2s}{2}.
\label{eq:lpt-machine-lb}
\end{equation}

\begin{equation}
C^*
\ge
2sn+\sum_{j:e_j<3s}(e_j-2s).
\label{eq:lpt-short-lb}
\end{equation}

\begin{theorem}[Exact LPT ratio for long jobs]
\label{thm:lpt-long-jobs}
For every integer \(s\ge2\), we obtain 
\[
\boxed{
\rho_{\mathrm{LPT}}^{\ge s}(s)
=
\frac{11s-3}{8s-2}.
}
\]
\end{theorem}

\begin{lemma}[Prefix--tail decomposition]
\label{lem:lpt-prefix-tail}
An LPT schedule has one of the following forms:

\begin{enumerate}
\item No short reset occurs. In this case, deleting every short job leaves the
      LPT makespan unchanged.
\item A first short reset occurs. The schedule can then be partitioned
      chronologically into a prefix \(P\) followed by a serial terminal part
      \(T\). The terminal part consists entirely of short jobs, except that it
      may begin with one long job forming a one-job block.
\end{enumerate}

Let \(L_P,E_P,n_P\) denote the length, execution volume, and number of jobs in
\(P\), and define
\[
I_P=2L_P-E_P.
\]
Then
\begin{equation}
I_P\le\frac{5s-2}{2}\,n_P.
\label{eq:lpt-prefix-idle}
\end{equation}
If \(E_T\) denotes the total execution length of the jobs in \(T\), then
\begin{equation}
C_{\mathrm{LPT}}=L_P+E_T.
\label{eq:lpt-prefix-tail-length}
\end{equation}
\end{lemma}

\begin{equation}
C_{\mathrm{LPT}}(I)=C_{\mathrm{LPT}}(I_L).
\label{eq:lpt-delete-short}
\end{equation}

\begin{theorem}[LPT bound for unrestricted instances]
\label{thm:lpt-unrestricted-upper}
For every integer \(s\ge2\),
\[
\boxed{
\rho_{\mathrm{LPT}}(s)
\le
\frac{13s-2}{8s}.
}
\]
\end{theorem}

\begin{proof}
For any \(\lambda\in[0,1]\), the convex combination of
\eqref{eq:lpt-machine-lb} and \eqref{eq:lpt-short-lb} gives
\begin{equation}
C^*
\ge
\lambda\frac{E+2s}{2}
+
(1-\lambda)
\left(
2sn+\sum_{j:e_j<3s}(e_j-2s)
\right).
\label{eq:lpt-combined-lb}
\end{equation}

The idle estimate in \eqref{eq:lpt-prefix-idle} can be covered whenever
\[
2s(1-\lambda)
\ge
\frac{\lambda(5s-2)}{4}.
\]
The largest admissible weight is
\begin{equation}
\lambda=\frac{8s}{13s-2}.
\label{eq:lpt-weight}
\end{equation}

If no short reset occurs, let \(I_L\) be the long-job subinstance. By
\eqref{eq:lpt-delete-short} and Theorem~\ref{thm:lpt-long-jobs}, we obtain 
\[
\begin{aligned}
C_{\mathrm{LPT}}(I)
&=C_{\mathrm{LPT}}(I_L)\\
&\le\frac{11s-3}{8s-2}C^*(I_L)\\
&\le\frac{11s-3}{8s-2}C^*(I),
\end{aligned}
\]
because deleting jobs cannot increase the optimal makespan. Finally,
\[
\frac{11s-3}{8s-2}<\frac{13s-2}{8s},
\qquad s\ge2.
\]

Suppose that a first short reset occurs and use the decomposition \(P\cup T\)
from Lemma~\ref{lem:lpt-prefix-tail}. We reserve the constant term
\(\lambda s\) in
\[
\lambda\frac{E+2s}{2}=\lambda\frac E2+\lambda s
\]
for the terminal part.

After discarding any additional short-job terms belonging to \(P\), the
contribution of the prefix to \eqref{eq:lpt-combined-lb} is at least
\[
\lambda\frac{E_P}{2}+2s(1-\lambda)n_P.
\]
Since \(E_P=2L_P-I_P\), we get 
\begin{equation}
\begin{aligned}
\lambda\frac{E_P}{2}+2s(1-\lambda)n_P
&=\lambda L_P-\frac{\lambda}{2}I_P
  +2s(1-\lambda)n_P\\
&\ge
\lambda L_P+
\left(
2s(1-\lambda)-\frac{\lambda(5s-2)}4
\right)n_P\\
&=\lambda L_P.
\end{aligned}
\label{eq:lpt-prefix-payment}
\end{equation}

If every job in \(T\) is short, its contribution to
\eqref{eq:lpt-short-lb} is \(E_T\). Including the reserved term \(\lambda s\),
the terminal contribution to \eqref{eq:lpt-combined-lb} is
\[
\lambda\frac{E_T}{2}+(1-\lambda)E_T+\lambda s.
\]
Its excess over \(\lambda E_T\) is
\[
\left(1-\frac{3\lambda}{2}\right)E_T+\lambda s.
\]
Since
\[
1-\frac{3\lambda}{2}=\frac{s-2}{13s-2}\ge0,
\]
the terminal contribution is at least \(\lambda E_T\).

It remains to consider the exceptional case in which \(T\) begins with one
long job. Let \(a\) be its execution length and \(b\) the length of the first
short job. Since the short job waits for a reset rather than overlapping the
long job,
\[
a-b<2s.
\]
Set
\[
d=a-2s.
\]
Because \(a\ge3s\), \(b\le3s-1\), and all quantities are integral,
\begin{equation}
s\le d\le3s-2,
\qquad
b\ge d+1.
\label{eq:lpt-tail-d}
\end{equation}
All terminal jobs after the first one are short. Their contribution to
\eqref{eq:lpt-short-lb}, together with the \(2s\) contribution of the first
job, is therefore
\[
E_T-d.
\]
Moreover,
\begin{equation}
E_T\ge a+b\ge2d+2s+1.
\label{eq:lpt-tail-volume}
\end{equation}

The terminal contribution is at least
\[
\lambda\frac{E_T}{2}
+(1-\lambda)(E_T-d)
+\lambda s.
\]
After subtracting \(\lambda E_T\) and applying
\eqref{eq:lpt-tail-volume}, the remaining amount is at least
\[
(1-2\lambda)d
+
\left(1-\frac{3\lambda}{2}\right)(2s+1)
+\lambda s.
\]
Because \(1-2\lambda<0\), this expression is minimised at
\(d=3s-2\). Substituting \eqref{eq:lpt-weight} gives
\[
\frac{(s-1)(s-2)}{13s-2}\ge0.
\]
Thus, this terminal part also contributes at least \(\lambda E_T\).

Combining the prefix and terminal estimates with
\eqref{eq:lpt-prefix-tail-length}, we obtain
\[
C^*
\ge
\lambda L_P+\lambda E_T
=
\lambda C_{\mathrm{LPT}}.
\]
Therefore,
\[
\frac{C_{\mathrm{LPT}}}{C^*}
\le
\frac1\lambda
=
\frac{13s-2}{8s}.
\]
\end{proof}

\end{document}
